\cgeentry{cge_model}{A standalone model, written whole in the user language}

\cgesec{Description}

\cmd{cge_model new} (from version 0.13.0) starts a standalone model in place
of the model in memory: a model written whole in the language of
\manref{cge_variable} and \manref{cge_equation}, without a SAM or the
built-in equations of Part~\ref{part:model}, as a GAMS program is. Its sets
are declared with \manref{cge_set}, its parameters with
\manref{cge_parameter}, as numbers, data or assignments that calibrate the
model, its variables with \manref{cge_variable}, at starting levels given by
expressions, and its equations with \manref{cge_equation}.
\manref{cge_solve} solves it. \cmd{cge_model describe} displays the
statistics of the model in memory, standalone or built from a SAM, as GAMS's
model statistics do. From version 0.17.0, \cmd{cge_model define} names a
model over a subset of the equations, as GAMS's \cmd{model} statement does,
which \manref{cge_solve}, \opt{}{model()}, solves on its own. From version
0.37.0, \cmd{cge_model roles} says which variables are the prices, the
consumption and the incomes of the model's groups, for \manref{cge_report}.

\cgesec{Quick start}

A new standalone model, with two goods, an endowment and a price for each
\begin{stcode}
cge_model new
cge_set g = g1 g2
cge_parameter e(g), values(g1:10 g2:5)
cge_variable P(g), value(1)
\end{stcode}
The statistics of the model in memory
\begin{stcode}
cge_model describe
\end{stcode}
A model of the equations \cmd{DEM} and \cmd{SUP}, solved on its own, and the
named models defined (from version 0.17.0)
\begin{stcode}
cge_model define markets = DEM SUP
cge_solve, name(mk) model(markets)
cge_model list
\end{stcode}

\cgesec{Syntax}

\begin{stsyntax}
\cmdb{cge_model new}\\
\cmdb{cge_model} \textbf{\opt{d}{escribe}} [\argu{model}]\\
\cmdb{cge_model define} \argu{model} \texttt{=} \argu{list}\\
\cmdb{cge_model drop} \argu{model}|\cmd{_all}\\
\cmdb{cge_model list}\\
\cmdb{cge_model roles}\texttt{,} \optarg{group}{s}{set} \optarg{good}{s}{set} \optarg{price}{s}{varname} \optarg{cons}{umption}{varname} [\optarg{inc}{ome}{varname}]\\
\cmdb{cge_model roles} [\texttt{,} \opt{}{list}|\opt{}{clear}]
\end{stsyntax}

\noindent A \argu{list} names user equations, variables defined by an
expression, inequalities and models defined before, or \cmd{all}, every
equation, separated by spaces or commas.

\cgesec{Remarks and examples}

A standalone model is built from the commands of the user language. The
letters \cmd{a c f h d i n} are not the indexes of model sets in it: its sets
take any names, with labels of their own or from the data in memory,
\cmd{cge_set ac, from(account)}; the index of a set with a digit after it,
such as \cmd{i2}, runs over the same set, and \cmd{cge_set j = i} defines a
set of the same elements, as GAMS's \cmd{alias}. Parameters are calibrated by
assignments, element by element, \cmd{cge_parameter share(c,a) = io(c,a) /
sum(c2, io(c2,a))}, which are evaluated once, as in GAMS. Variables start at
the levels of \opt{}{value()}, an expression that may use the parameters,
and exogenous variables are \opt{}{fixed}; \manref{cge_closure} exchanges
fixed and free variables with \opt{}{swap()}. The starting levels need not
solve the model: \manref{cge_solve} solves it from them, and
\manref{cge_results} compares each solution with them.

\cgesubsec{Example 1: an exchange economy}

Two consumers trade two goods, with Cobb-Douglas utility and endowments
$e_{hg}$; good 1 is the numeraire:
\stlog{model_new}
The model has nine equations in nine free variables:
\stlog{model_describe}
With $\alpha_{hg}$ the budget shares, the market for good 1 gives the price
of good 2, $P_2 = \sum_h \alpha_{h2}\,e_{h1} / \sum_h \alpha_{h1}\,e_{h2}$,
here $8.2/5.4 = 1.5185$, which the solution reproduces; \cmd{WALRAS}, the
excess supply of good 1, is zero, as Walras's law requires:
\stlog{model_solve}

\cgesubsec{Example 2: a scenario}

The endowments are parameters, which \manref{cge_shock} changes:
\stlog{model_shock}
The market equations use the endowments themselves. Had they used a
parameter assigned from them, \cmd{cge_parameter etot(g) = sum(h,
e(h,g))}, the shock would not reach it: an assignment is evaluated once.

\cgesubsec{Example 3: named models}

From version 0.17.0, a named model is a subset of the equations, as in GAMS:
\cmd{model markets / DEM, SUP /}. Two markets with supply $2P^2$ and demand
$a/\sqrt{P}$, and the value of their sales as a defined variable; the
markets alone determine the prices and quantities, whose closed form is $P =
(a/2)^{2/5}$:
\stlog{model_named}
Solved on its own, a model takes its equations and the free variables that
appear in them; the other variables keep their levels, here \cmd{TOT} its
starting level of 2. \opt{}{from()} starts a solve from the levels of a
stored simulation, as GAMS's levels carry over from one solve to the next,
and \opt{}{hold()} holds variables at their levels for one solve, as GAMS's
\cmd{P.fx = P.l}: the total solved from the markets' solution, with the prices
and quantities held, is that of the whole model.

\cgesubsec{Example 4: the roles of the variables}

\manref{cge_report} knows the households, prices and consumption of a model
built from a SAM. Those of a standalone model are its own variables, and
\cmd{cge_model roles} names them: the set of groups (households, or any
others), the set of goods, the variable of the prices the groups pay,
indexed by the goods, that of the quantities each group buys, indexed by the
goods and the groups in either order, and, optionally, that of each group's
income:
\stlog{model_roles}
\manref{cge_report} then gives each group's expenditure, price indexes and
welfare (\manref{cge_report}, Example 5). The roles are saved with the model.

\cgesubsec{Example 5: the standard model}

The example \cmd{examples/standard_model.do} of the package repository
writes the standard model of Part~\ref{part:model} for the synthetic SAM as a
standalone model: 6 sets, 48 parameters calibrated from the SAM by
assignments, 42 blocks of variables and 35 blocks of equations, 93 equations
in all. Its solutions agree with those of the built-in model to about
$10^{-12}$ at the base and in four scenarios, and GAMS solves it to the same
solution (\chapref{chap:validation}).

\cgesec{Stored results}

\cmd{cge_model new} stores \cmd{r(standalone)}, 1. \cmd{cge_model describe}
stores the following in \cmd{r()}:
\begin{storedres}
\resgroup{Scalars}
\res{r(standalone)} & 1 for a standalone model, 0 for a model from a SAM\\
\res{r(N_eq_blocks)} & blocks of equations\\
\res{r(N_equations)} & single equations\\
\res{r(N_var_blocks)} & blocks of variables\\
\res{r(N_variables)} & single variables\\
\res{r(N_fixed)} & fixed variables\\
\res{r(N_free)} & free variables\\
\res{r(N_nonzeros)} & non-zero entries of the Jacobian at the current point\\
\res{r(N_sets)} & user sets\\
\res{r(N_parameters)} & user parameters\\
\res{r(N_constraints)} & constraints\\
\res{r(square)} & 1 if the system is square\\
\res{r(N_reduced)} & equations of the reduced system (a model from a SAM)\\
\res{r(N_models)} & named models defined (from version 0.17.0)\\
\resgroup{Macros}
\res{r(model)} & the named model described, if any\\
\end{storedres}
\cmd{cge_model define} stores \cmd{r(name)}, \cmd{r(equations)}, the equations of
the model after models and \cmd{all} are expanded, and \cmd{r(N_eq_blocks)};
\cmd{cge_model list} stores \cmd{r(models)} and \cmd{r(N_models)}.
\cmd{cge_model roles} stores \cmd{r(groups)}, \cmd{r(goods)},
\cmd{r(prices)}, \cmd{r(consumption)} and \cmd{r(income)}, empty without
roles.

\cgesec{Methods and formulas}

A standalone model has only user variables, in the order they were added,
and two blocks of equations: the user equations and the constraints of their
inequalities (\secref{sec:standalone}). Every equation element is solved,
with the exact Jacobian of their postfix rows; there is no reduced system.
Its closure is that of the variables: those added \opt{}{fixed}, then the
swaps. \manref{cge_gams} writes it as a GAMS program of its own: its labels
as the set \cmd{AC}, its parameters, variables and equations, and a driver
that solves the starting levels, with the gaps of the equations there, and
then the scenario.

A named model (from version 0.17.0) is kept as its list, expanded when it is
solved, so that it follows later equations. Its system has the elements of
its user equations and defined variables, the constraints of its
inequalities and the free variables that appear in them, found from the
pattern of the Jacobian; the other free variables are fixed at their levels
at the start of the solve. Solved as a system, it must be square; with
\opt{min}{imize()} or \opt{max}{imize()} it is a nonlinear programme over its
variables. \manref{cge_gams} exports a simulation of a named model with
\cmd{model cgesub / UE}$k$\cmd{, ... /} and the variables outside it fixed at
their levels (\cmd{.fx}).

\cgesec{Also see}

\begin{alsosee}
\seeentry{cge_set}{User sets of elements, for indexes and sums}
\seeentry{cge_parameter}{User parameters: values, assignments and data}
\seeentry{cge_variable}{User variables, defined by an expression or determined by user equations}
\seeentry{cge_equation}{User equations: new relations among the model's variables}
\seechapter{chap:user}{User variables and equations}
\end{alsosee}
